\chapter{Ipparchia dove sei?}

«Ippa! Ippa! Ipparchia!»\\
La chiama forte, ma lei non torna.

Si siede sul muretto di cinta di un edificio in rovina che non ha mai esplorato, per via di un odore che lo preoccupa: una punta sottile di solvente, insidiosa, che ogni volta gli fa prudere le narici come se volesse raggiungergli il cervello.

Non sa cosa sia stato abbandonato oltre quel muretto e non ha mai cercato di scoprirlo.

Appoggia il cartone sulle ginocchia e aspetta.

Una volta, poco dopo aver perso la vista, gli dissero di attendere il medico: sarebbe arrivato a momenti. Giovanni attese due ore, seduto buono, finché un’infermiera si accorse di lui e lo accompagnò alla visita.

Da allora aveva imparato che “a momenti” può voler dire molte cose. Così, quando deve aspettare, recita i suoi autori preferiti. Spesso Schopenhauer, \textit{L’arte di ottenere ragione}.

«Stratagemma numero uno...»

Lo ripete tutto e, alla fine, sa che è passato circa un minuto. Se fosse stato un matematico avrebbe contato: uno, due, tre. Ma la matematica non era la sua passione.

Arriva allo stratagemma trentotto: è passata un’ora, e Ipparchia non è ancora tornata.

Dietro di lui, un luogo ignoto. Alla sua destra, la strada da cui è arrivato; a sinistra, la strada verso casa. Davanti, una via semi-inesplorata costeggiata da una fascia ecologica, il  margine vivo in cui Ipparchia deve essersi infilata seguendo il fagiano. \'E la direzione più probabile.
Si incammina, continuando a chiamarla.

La strada diventa bianca: l’asfalto finisce e comincia la terra battuta. La fascia ripariale, però, non si interrompe, continua  come un corridoio di odori, il percorso più probabile lungo cui Ipparchia può aver proseguito l’inseguimento.

Quando era scattata dietro al fagiano erano circa le otto di sera; ora è quasi buio. Per Giovanni cambia poco, ma i cani usano anche la vista.

Arriva a un incrocio: la strada bianca curva a destra, la fascia si interrompe, mentre davanti a lui l’asfalto ricomincia.

«Ippa! Dove sei andata?»

Giovanni tiene la destra e si addentra in campagna. Gli alberi continuano a costeggiare la strada, ne conta nove, poi all’improvviso l’asfalto ritorna.

\par\medskip
\begin{center}
  \includegraphics[width=0.9\linewidth]{mappa con casa laura.png}
\end{center}

Attende qualche minuto. Sente passare un veicolo. La strada che ha incontrato taglia quella bianca: deve essere una via principale.

Sta per tornare indietro, quando dall’altra parte arriva un latrato familiare.

«Ippa!»

Attraversa senza esitazione.

«Ippa!»

Le va incontro. Lei rimane ferma. Quando la raggiunge, capisce: il guinzaglio di corda è rimasto incastrato in qualche anfratto. Giovanni lo libera facilmente.

Ipparchia gli lecca la mano.

Ma le lingue sono due.

Un altro cane è accanto a lei.

«C’è un altro cane qui...»

Chiama più volte:

«C’è un cane qui! È di qualcuno questo cane?»

Nessuna risposta.\\

Aspetta ancora un poco. Non abbastanza da arrivare di nuovo allo stratagemma trentotto.\\

«Andiamo, Ippa.»

Giovanni si rivolge solo a lei, ma la coppia ha fatto un acquisto: qualcuno li segue scodinzolando.

Raggiungono di nuovo la strada bianca. Giovanni si inginocchia e accarezza il nuovo arrivato.

«E tu? Ci segui?»

Si gratta la testa.

«In ogni caso, non saprei come altro aiutarti, per ora. Va bene: andiamo a casa, poi vedremo.»

Ripercorre la strada dell’andata, passo dopo passo, rumore dopo rumore, odore dopo odore.\\
«Coraggio, è tardi, ma si va a cena.»

Giovanni regge il pacco per le ultime decine di metri che lo separano dal loro rifugio.\
L’odore delle rose dell’ultima abitazione si mescola alla lavanda spontanea che cresce nel parco abbandonato dell’ex convitto.\
Avanza tra spine di more e ciuffi di rosmarino, che pungono le gambe e coprono l’odore di Ippa e dell’altro cane. Un dedalo verde che è la sua protezione segreta.

Le braccia gli pesano, ma ormai ce l’ha fatta.

Tra le foglie filtra un filo di luce. Viene da una candela accesa nel convitto, passa tra le assi che chiudono una finestra senza vetro e arriva fino a lui. Poche cose come una luce accesa sanno dire “casa” senza parole.

Ma Giovanni non la può vedere.

Come non può vedere altre cose intorno a lui, ma ci sente benissimo: un rametto secco si spezza a pochi metri da lui. Non è un gatto nè un riccio, troppo leggeri.\
«Ehi, c’è qualcuno?»\
Silenzio.

Nessuna risposta.

Fa un passo verso la scala antincendio, ma alle sue spalle una voce lo ferma:\
«Appoggia il pacco, amico.»

Giovanni resta immobile.\
«Mi hai trovato.»\
Lo dice lentamente, come una battuta preparata da anni.

Poi posa il pacco sul primo pianerottolo e lascia scorrere la mano verso il fianco, quasi senza muoversi.

«Io non lo farei…»

La voce è leggermente coperta   dal nuovo amico di Ipparchia che ringhia sommessamente. Ha il pelo irto, le zampe piantate a terra, lo sguardo fisso verso il buio alle sue spalle.

«Tu non sei me» dice Giovanni. «Perché tu sei…»

Un passo nel buio. Qualcosa striscia tra le foglie secche, appena oltre il cerchio di luce.

Giovanni stringe le mascelle.

Ipparchia si sdraia, l’altro cane ringhia più forte.

Poi la voce, ferma e bassa:

«…un uomo morto.»

Un colpo secco, come legno che batte sul metallo, rompe il silenzio.

